\section{What is not verified; open questions}\label{sec:limitations}

\paragraph{Not formally verified.}
The following are proved on paper only.
\begin{itemize}
\item Proposition~\ref{prop:forced}, including the reduction from line arrangements to wiring diagrams.
\item Lemma~\ref{lem:percycle}, the sweep argument.
\item Lemma~\ref{lem:A}, which relies on the Matsumoto--Tits theorem~\cite{Matsumoto1964}.
\item Hence Theorem~\ref{thm:parity}, Corollary~\ref{cor:main}, and the upper bounds $M(7)\le18$ and $M(9)\le33$.
\end{itemize}

An attempt to certify Lemma~\ref{lem:A} for $n=7$ in Lean, by a kernel-evaluated breadth-first search over
the sweep states of all $5040$ permutations, ran out of kernel memory. We chose not to fall back on
\texttt{native\_decide}.

\paragraph{Computational evidence.} This is not a proof. The sweep process in the proof of
Lemma~\ref{lem:percycle} defines a finite-state dynamic program. We validated it against direct component
counts of $G$ on random reduced words for $n=3,5,\dots,11$, and then ran it over \emph{all} reduced
words of $w_0$ for $n\le 9$. For $n=7$ and $n=9$ the graph $G$ is never a single cycle; it has $2$ or at
least $3$ components. In every case $\#A$ had the same parity for all words, as Lemma~\ref{lem:A} predicts.

On $2000$ random reduced words for each odd $n\le 19$, the parity of the number of components always
equalled $(n^2-1)/8\bmod 2$. On $2500$ random words (500 for each $n=5,7,9,11,13$), every cycle
contained an odd number of $A$-crossings. Random real arrangements of $7$ lines produced $2$, $4$ or
$6$ components, never an odd number.

\paragraph{Open questions.}
\begin{enumerate}
\item Is $M(n)=\binom n2$ for every $n\equiv\pm3\pmod 8$? We confirmed this for $n=3,5,11,13$. Brief
  searches for $n=19$ and $n=21$ (four runs of $10^6$ annealing steps each) reached $3$ components but
  not $1$.
\item Is $M(n)=\binom n2-3$ for every $n\equiv\pm1\pmod 8$? We confirmed this for $n=7$ and $n=9$.
\item What is $M(n)$ for even $n$? The closed-trail bound is $n(n-2)/2$, and DSM's $24$-gon on $8$
  lines attains it for $n=8$. We have not studied even $n$.
\item Is there a proof of Lemma~\ref{lem:A} that avoids the Matsumoto--Tits theorem, for example via an
  explicit potential function on permutations? A linear-algebra search over $\mathrm{GF}(2)$ suggests
  that such potentials exist.
\end{enumerate}
